\section{Design Choices}

\subsection{A Finitary Oracle Model}

The project deliberately does not start from a high-level semantic type of
oracle functionals such as \(Set\ \N \to \N \to. \N\).  Such a choice would
require substantial infrastructure before the priority construction could
begin: an oracle-machine model, continuity or compactness theorems,
extraction of finite uses from semantic computations, and a connection
between semantic functionals and executable indices.

Instead, finite execution is primary.  The evaluator runs a concrete
program for bounded fuel against a partial oracle \(\N\to Option\ Bool\).
Infinite computation is later defined by existence of a finite successful
run against a total oracle.  This representation puts exactly the finite
information used by a priority construction into the main semantic object.

\subsection{Finite Sets Versus Boolean Prefixes}

The construction represents evolving c.e. approximations as monotone finite
sets:

\begin{lstlisting}
def StageMono (F : Nat -> Finset Nat) : Prop :=
  forall s, F s subset F (s + 1)

def limitSet (F : Nat -> Finset Nat) : Set Nat :=
  {n | exists s, n in F s}
\end{lstlisting}

The actual syntax uses Lean's \texttt{\(\subseteq\)} and \texttt{\(\in\)}
notation.  Boolean lists appear separately as snapshots.  This separation
is one of the important design decisions: a stage set can enumerate new
positive information at any later stage, while a snapshot merely supplies a
finite window of answers to one bounded oracle computation.

\subsection{Hybrid Use of Mathlib}

The development separates project-local oracle reasoning from ordinary
computability closure proofs.  Project-local files define oracle syntax,
bounded evaluation, infinite semantics, reducibility, the construction,
invariants, finite injury, and satisfaction of requirements.  Mathlib's
computability library enters through \texttt{MathlibBridge.lean},
\texttt{RunPrimrec.lean}, and the c.e. proof in \texttt{CE.lean}.  This
boundary allows the priority proof to reason about explicit finite uses
while still reusing Mathlib's closure theorems for primitive-recursive and
partial-recursive functions \cite{carneiro2019,mathlibPartrec}.

\subsection{Local Reducibility Vocabulary}

The central reducibility definition is:

\begin{lstlisting}
def TuringReducible (A B : Set Nat) : Prop :=
  exists c : OracleCode,
    forall x : Nat,
      Computes c (charFun B) x (natOfBool (charFun A x))
\end{lstlisting}

Thus \(A\leq_T B\) means that a project-local oracle program computes the
characteristic function of \(A\) using the total characteristic function of
\(B\) as oracle.  The development proves reflexivity
\texttt{TuringReducible.refl}, that computable sets reduce to every oracle
via \texttt{ComputableSet.turingReducible}, and an effective bijection
between natural numbers and \texttt{OracleCode}.  It does not yet prove
composition or transitivity of \texttt{TuringReducible}; that limitation is
discussed in Section~\ref{sec:limitations}.

